%============================================================
% CH1 · P4 · PROBLEM 2 — OPEN THE LOCK (LeetCode 752)
%============================================================
\subsection{Problem 1.4.2: Open the Lock}
\problemlabel{1.4.2}{ssec:ch1-open-lock}

\problemcard{1.4.2}{Open the Lock}{LeetCode 752}{https://leetcode.com/problems/open-the-lock/}{\exmedium\ \lc{Medium}}{Codes as vertices of a 4-dimensional torus, forbidden vertices, BFS; bidirectional BFS}

\subsubsection*{1. Problem statement}

A lock has 4 circular wheels, each showing a digit \texttt{0}--\texttt{9}. One move turns \textbf{one} wheel by \textbf{one} slot, up or down, with wrap-around (\texttt{9} $\to$ \texttt{0} and \texttt{0} $\to$ \texttt{9}). The lock starts at \texttt{"0000"}. A list of \textbf{dead ends} is given: if the lock ever displays one of them, the wheels jam and it can never be opened. Return the minimum number of moves to reach the code \texttt{target}, or $-1$ if it is impossible.

\textbf{Constraints:} $1 \le |\texttt{deadends}| \le 500$; all codes have 4 digits; \texttt{target} is not a dead end.

\begin{table}[H]
\centering\small
\begin{tabular}{@{}p{5.6cm}ccp{5.6cm}@{}}
\toprule
\textbf{deadends} & \textbf{target} & \textbf{Output} & \textbf{Comment} \\
\midrule
\texttt{0201 0101 0102 1212 2002} & \texttt{0202} & 6 & direct route (4 moves) is blocked \\
\texttt{8888} & \texttt{0009} & 1 & turn the last wheel down: $0 \to 9$ \\
\texttt{8887 8889 8878 8898 8788 8988 7888 9888} & \texttt{8888} & $-1$ & all 8 neighbours of the target are dead \\
\bottomrule
\end{tabular}
\end{table}

\begin{diagrambox}
\begin{tikzpicture}
    % lock body
    \draw[rounded corners=5pt, draw=deepblue, fill=theorybg, line width=1pt] (-0.4,-0.7) rectangle (4.0,0.9);
    \draw[line width=3pt, draw=midgray, rounded corners=8pt] (0.6,0.9) -- (0.6,1.7) -- (3.0,1.7) -- (3.0,0.9);
    \foreach \d [count=\i from 0] in {0,0,0,0}
        \node[gcell, minimum width=8mm, minimum height=10mm, font=\Large\ttfamily\bfseries] at (\i*0.9+0.4,0.1) {\d};
    \node[font=\scriptsize, text=darkgray] at (1.75,-1.05) {start \texttt{0000}, target \texttt{0202}};
    % one wheel as a ring
    \begin{scope}[xshift=7.3cm, yshift=0.2cm]
        \foreach \k in {0,...,9} {
            \pgfmathsetmacro{\ang}{90-36*\k}
            \node[gv, minimum size=6mm, font=\scriptsize\bfseries] (r\k) at (\ang:1.35) {\k};
        }
        \foreach \k in {0,...,9} {
            \pgfmathtruncatemacro{\nx}{mod(\k+1,10)}
            \draw[ge] (r\k) -- (r\nx);
        }
        \draw[gpath] (r9) -- (r0);
        \node[font=\scriptsize, text=accentorange, anchor=west] at (1.75,1.45) {\texttt{9} $\leftrightarrow$ \texttt{0}: one move};
        \node[note] at (0,-2.0) {each wheel is a cycle of length 10};
    \end{scope}
\end{tikzpicture}
\figcap{The lock and one of its wheels. Turning a wheel moves its digit to a neighbour on the cycle; \texttt{0} and \texttt{9} are neighbours too.}
\end{diagrambox}

\subsubsection*{2. The graph model}

\begin{modelbox}[Graph Model --- Open the Lock]
\renewcommand{\arraystretch}{1.35}
\begin{tabularx}{\linewidth}{@{}>{\bfseries}l X@{}}
Vertices & all $10^4$ codes \texttt{0000}..\texttt{9999}, encoded as the integer $0..9999$; dead ends are \textbf{deleted} vertices \\
Edges & undirected: two codes are adjacent if they differ in exactly one wheel by $\pm1 \pmod{10}$; every code has exactly 8 neighbours, generated on the fly \\
Weights & none --- each turn costs 1 \\
Question & shortest path from \texttt{0000} to \texttt{target} avoiding deleted vertices ($-1$ if disconnected) \\
Size & $|V| = 10^4$, \ $|E| = 4\cdot10^4$ (8 per vertex, each counted twice): tiny \\
\end{tabularx}
\end{modelbox}

\begin{diagrambox}
\begin{tikzpicture}
    \node[gvsrc, minimum size=11mm, font=\small\ttfamily\bfseries] (c) at (0,0) {0000};
    \foreach \lab/\ang/\w in {1000/90/1, 9000/45/1, 0100/0/2, 0900/-45/2, 0010/-90/3, 0090/-135/3, 0001/180/4, 0009/135/4} {
        \node[gvseen, minimum size=11mm, font=\scriptsize\ttfamily\bfseries] (n\lab) at (\ang:2.3) {\lab};
        \draw[ge] (c) -- (n\lab);
    }
    \node[note, anchor=west, text width=6.2cm, align=left] at (3.4,0.4) {
        \textbf{8 neighbours} of every code: wheel $1..4$, turned up or down.\\[3pt]
        Seen as a whole, the graph is a 4-dimensional grid of side 10 whose rows wrap around (a \emph{torus}): Pattern~3's grid, but in four dimensions, with dead ends playing the role of walls.};
\end{tikzpicture}
\figcap{The neighbourhood of \texttt{0000}. All 8 neighbours are at distance 1 (unless they are dead ends).}
\end{diagrambox}

\subsubsection*{3. Key observations}

\begin{itemize}[leftmargin=*]
    \item \textbf{Fewest moves, all moves equal} $\Rightarrow$ BFS from \texttt{0000}.
    \item \textbf{Without dead ends the answer is a formula:} each wheel is independent, so the distance is $\sum_{i} \min(d_i, 10 - d_i)$ over the target digits $d_i$. For \texttt{0202} that is $2 + 2 = 4$. Dead ends force \textbf{detours} --- sample 1 needs 6 --- and only a search sees them.
    \item \textbf{Dead ends are deleted vertices}, not just ``bad answers'': a path may not pass \emph{through} them. Never push them into the queue.
    \item \textbf{If \texttt{0000} itself is a dead end}, the lock is jammed from the start: answer $-1$ (unless the target were \texttt{0000}, which the statement excludes by forbidding a dead target).
    \item \textbf{The state space is fixed and small} ($10^4$ codes): no window proof is needed, unlike Problem~1.4.1.
\end{itemize}

\begin{diagrambox}
\begin{tikzpicture}
    \node[gvdst, minimum size=11mm, font=\small\ttfamily\bfseries] (c) at (0,0) {0202};
    \foreach \lab/\ang/\st/\d in {1202/90/gvseen/5, 9202/45/gvseen/5, 0302/0/gvseen/5, 0102/-45/gvdead/{dead}, 0212/-90/gvseen/5, 0292/-135/gvseen/5, 0201/180/gvdead/{dead}, 0203/135/gvseen/5} {
        \node[\st, minimum size=11mm, font=\scriptsize\ttfamily\bfseries] (n\lab) at (\ang:2.4) {\lab};
        \node[gdist] at ($(n\lab)+(\ang:0.85)$) {\d};
    }
    \foreach \lab in {1202, 9202, 0302, 0212, 0292, 0203} \draw[ge] (c) -- (n\lab);
    \foreach \lab in {0102, 0201} \draw[gfaded, dashed] (c) -- (n\lab);
    \node[note, anchor=west, text width=6cm, align=left] at (3.5,0.2) {
        The two neighbours on a 4-move route, \texttt{0201} and \texttt{0102}, are dead.
        Every live neighbour is at distance 5 (green), so the answer is $5 + 1 = 6$.};
\end{tikzpicture}
\figcap{Sample 1 seen from the target. Distances in green were computed by a full BFS from \texttt{0000}.}
\end{diagrambox}

\subsubsection*{4. Approaches}

\begin{bruteforcebox}
\textbf{Idea: try all move sequences}, shortest first (iterative deepening): for $d = 0, 1, 2, \dots$ enumerate every sequence of $d$ turns that avoids dead ends.

\begin{lstlisting}[style=cpp]
#include <bits/stdc++.h>
using namespace std;

class Solution {
    set<string> dead;
    string target;
    bool reach(string code, int k) {                 // target within k more moves?
        if (code == target) return true;
        if (k == 0) return false;
        for (int w = 0; w < 4; w++)
            for (int step : {1, 9}) {                // +1 or -1 (mod 10)
                string nxt = code;
                nxt[w] = char('0' + (code[w] - '0' + step) % 10);
                if (!dead.count(nxt) && reach(nxt, k - 1)) return true;
            }
        return false;
    }
public:
    int openLock(vector<string>& deadends, string target) {
        dead = set<string>(deadends.begin(), deadends.end());
        this->target = target;
        if (dead.count("0000")) return -1;
        for (int d = 0; d <= 20; d++)                // no answer exceeds 20 (see trace)
            if (reach("0000", d)) return d;
        return -1;                                   // may take forever to get here
    }
};
\end{lstlisting}

\begin{complexitybox}
\textbf{Time:} $\Theta(8^{d})$ sequences for answer $d$: sample 1 ($d=6$) already makes 51135 recursive calls (depths $0..6$ together), an answer of 20 would need $8^{20} \approx 10^{18}$, and an impossible case runs all depths up to 20. The same code is reached by a huge number of different sequences. \textbf{Verdict:} \tle{}.
\end{complexitybox}
\end{bruteforcebox}

\begin{algorithmbox}
\textbf{Idea: BFS over the $10^4$ codes}, each code an integer, \texttt{dist} a flat array, dead ends pre-marked as forbidden. Turning wheel $w$ (place value $p = 1000, 100, 10, 1$) from digit $g$ to digit $g'$ changes the integer by $(g' - g)\cdot p$.

\begin{lstlisting}[style=cpp]
#include <bits/stdc++.h>
using namespace std;

class Solution {
public:
    int openLock(vector<string>& deadends, string target) {
        const int P[4] = {1000, 100, 10, 1};           // place value of each wheel
        vector<char> dead(10000, 0);
        for (const string& s : deadends) dead[stoi(s)] = 1;
        int goal = stoi(target);
        if (dead[0]) return -1;                         // jammed from the start
        vector<int> dist(10000, -1);
        queue<int> q;
        dist[0] = 0;
        q.push(0);
        while (!q.empty()) {
            int u = q.front(); q.pop();
            if (u == goal) return dist[u];
            for (int w = 0; w < 4; w++) {
                int g = u / P[w] % 10;                  // digit on wheel w
                for (int step : {1, 9}) {               // up, down (9 = -1 mod 10)
                    int v = u + ((g + step) % 10 - g) * P[w];
                    if (dead[v] || dist[v] != -1) continue;
                    dist[v] = dist[u] + 1;
                    q.push(v);
                }
            }
        }
        return -1;                                      // target unreachable
    }
};
\end{lstlisting}

\begin{complexitybox}
\textbf{Time:} $\bigO{10^4 \cdot 8}$ plus reading the dead ends --- about $10^5$ steps whatever the input. \textbf{Space:} $\bigO{10^4}$. \textbf{Verdict:} \ac{}.
\end{complexitybox}
\end{algorithmbox}

\begin{invariantbox}[Why it is correct]
The graph has exactly the live codes as vertices and the legal single turns as edges, so a sequence of moves that never shows a dead end is exactly a path from \texttt{0000} in this graph. Dead codes are never pushed, so no path through them is ever produced. BFS with unit costs sets \texttt{dist[v]} to the true minimum number of moves the first time $v$ is discovered (Pattern~4 invariant, \S\ref{ssec:ch1-p4-theory}); when \texttt{goal} is popped its distance is final. If the queue empties first, the goal is not in the component of \texttt{0000}: $-1$.
\end{invariantbox}

\begin{alternativebox}[Approach 3: Bidirectional BFS]
\textbf{Idea: search from both ends and meet in the middle.} The graph is undirected and the goal is a single known code, so we may grow a BFS ball around \texttt{0000} \emph{and} one around \texttt{target}, always expanding one full layer of the \textbf{smaller} frontier. As soon as an edge connects the two balls, the shortest path is found.

\begin{diagrambox}
\begin{tikzpicture}
    \fill[primaryblue!12] (0,0) circle (1.5);
    \fill[successgreen!12] (3.2,0) circle (1.5);
    \draw[primaryblue!60, dashed] (0,0) circle (1.5);
    \draw[successgreen!60, dashed] (3.2,0) circle (1.5);
    \node[gvsrc, minimum size=8mm, font=\scriptsize\ttfamily\bfseries] at (0,0) {0000};
    \node[gvdst, minimum size=8mm, font=\scriptsize\ttfamily\bfseries] at (3.2,0) {0202};
    \node[gv, minimum size=4mm] (x) at (1.45,0.3) {};
    \node[gv, minimum size=4mm] (y) at (1.8,0.3) {};
    \draw[gpath] (x) -- (y);
    \node[font=\scriptsize, text=primaryblue] at (0,-1.8) {radius 3 from start};
    \node[font=\scriptsize, text=successgreen!70!black] at (3.2,-1.8) {radius 2 from target};
    \draw[primaryblue!40, dashed] (9.0,0) circle (2.4);
    \fill[primaryblue!8] (9.0,0) circle (2.4);
    \node[gvsrc, minimum size=8mm, font=\scriptsize\ttfamily\bfseries] at (7.8,0) {0000};
    \node[gvdst, minimum size=8mm, font=\scriptsize\ttfamily\bfseries] at (10.2,0) {0202};
    \node[font=\scriptsize, text=primaryblue] at (9.0,-2.7) {one-sided BFS: radius 6};
\end{tikzpicture}
\figcap{Two balls of radius about $d/2$ instead of one ball of radius $d$. In a graph where the number of states grows quickly with the radius, that is a huge saving.}
\end{diagrambox}

\begin{lstlisting}[style=cpp]
#include <bits/stdc++.h>
using namespace std;

class Solution {
public:
    int openLock(vector<string>& deadends, string target) {
        const int P[4] = {1000, 100, 10, 1};
        vector<char> dead(10000, 0);
        for (const string& s : deadends) dead[stoi(s)] = 1;
        int goal = stoi(target);
        if (dead[0]) return -1;
        if (goal == 0) return 0;
        vector<int> da(10000, -1), db(10000, -1);       // distances from start / from goal
        vector<int> fa = {0}, fb = {goal};              // current frontiers (one layer each)
        da[0] = 0;
        db[goal] = 0;
        while (!fa.empty() && !fb.empty()) {
            bool fromStart = fa.size() <= fb.size();    // expand the smaller side
            vector<int>& front = fromStart ? fa : fb;
            vector<int>& mine = fromStart ? da : db;
            vector<int>& other = fromStart ? db : da;
            vector<int> next;
            int best = INT_MAX;
            for (int u : front)
                for (int w = 0; w < 4; w++) {
                    int g = u / P[w] % 10;
                    for (int step : {1, 9}) {
                        int v = u + ((g + step) % 10 - g) * P[w];
                        if (dead[v]) continue;
                        if (other[v] != -1) best = min(best, mine[u] + 1 + other[v]);
                        if (mine[v] == -1) { mine[v] = mine[u] + 1; next.push_back(v); }
                    }
                }
            if (best != INT_MAX) return best;           // the two balls touched
            front.swap(next);
        }
        return -1;                                      // one side got stuck: no path
    }
};
\end{lstlisting}

\textbf{Correctness.} Every layer is expanded \emph{completely} before checking, and the minimum over all touching edges of that layer is returned, so the first contact gives the shortest total. If one frontier becomes empty, its whole component has been explored without touching the other side: $-1$ (this is how the walled-in target of sample 3 is detected after a single layer).

\begin{complexitybox}
\textbf{Time:} at most $\bigO{10^4 \cdot 8}$ like BFS, often much less. On sample 1 one-sided BFS discovers 1562 codes before popping the target; bidirectional BFS discovers 245 (counting both sides). On \texttt{5555} with no relevant dead ends (distance 20, the maximum) there is no gain: both balls together cover the whole lock. \textbf{Verdict:} \ac{}; the technique pays off in huge state spaces such as sliding puzzles and word ladders (Chapter 3).
\end{complexitybox}
\end{alternativebox}

\begin{tricksbox}[Approach 4: Implementation Refinements]
\begin{itemize}[leftmargin=*]
    \item \textbf{Integers, not strings.} \texttt{unordered\_set<string>} and \texttt{string} copies work (Template~3 of the theory) but are several times slower; the integer id with \texttt{P[w]} arithmetic needs no allocation at all.
    \item \textbf{Down $=$ up by 9.} Writing $(g + 9) \bmod 10$ for ``down'' avoids negative remainders: in C++, $(0 - 1) \bmod 10 = -1$, not 9.
    \item \textbf{Pre-mark dead ends as visited.} Setting \texttt{dist[dead] = -2} merges the two tests into one; the code above keeps them separate for clarity.
    \item \textbf{Quick lower bound.} $\sum_i \min(d_i, 10-d_i)$ is a lower bound on the answer; if BFS reaches it, no detour was needed. Using it to guide the search is the idea of A* search.
\end{itemize}
\end{tricksbox}

\subsubsection*{5. Trace}

The lock graph is 4-dimensional, but sample 1 lives in a 2-dimensional \emph{slice}: the codes \texttt{0a0b} (wheels 1 and 3 stay at 0). In this slice the graph is an ordinary grid and the dead ends are walls --- a Pattern~3 picture.

\begin{tracebox}[{Sample 1 on the slice \texttt{0a0b}: distances from \texttt{0000}}]
\centering
\begin{tikzpicture}[x=0.95cm, y=0.95cm]
    % values row a (0..3), col b (0..4); X = dead
    \foreach \a/\row in {0/{0,1,2,3,4}, 1/{1,-,-,4,5}, 2/{2,-,6,5,6}, 3/{3,4,5,6,7}} {
        \foreach \v [count=\b from 0] in \row
            \node[gcell, minimum size=9mm] at (\b,\a) {\v};
    }
    % dead ends, path cells, start, target
    \foreach \b/\a in {1/1, 2/1, 1/2}
        \node[gwall, minimum size=9mm, font=\scriptsize\bfseries] at (\b,\a) {dead};
    \foreach \b/\a/\v in {0/1/1, 0/2/2, 0/3/3, 1/3/4, 2/3/5}
        \node[gcell, minimum size=9mm, fill=pathfill, font=\small\bfseries] at (\b,\a) {\v};
    \node[gcell, minimum size=9mm, fill=basefill, draw=successgreen, line width=1.2pt] at (0,0) {\textbf{0}};
    \node[gcell, minimum size=9mm, fill=answerfill, draw=accentorange, line width=1.2pt] at (2,2) {\textbf{6}};
    \foreach \b in {0,...,4} \node[font=\scriptsize, text=darkgray] at (\b,-0.75) {$b{=}\b$};
    \foreach \a in {0,...,3} \node[font=\scriptsize, text=darkgray] at (-0.95,\a) {$a{=}\a$};
    \node[note, anchor=west, text width=5.6cm, align=left] at (5.0,1.5) {
        cell $(a,b)$ = code \texttt{0a0b}, number = BFS distance.\\[3pt]
        The direct routes into \texttt{0202} pass \texttt{0201} or \texttt{0102} --- both dead.
        One shortest detour (shaded cells): \texttt{0000}$\to$\texttt{0100}$\to$\texttt{0200}$\to$\texttt{0300}$\to$\texttt{0301}$\to$\texttt{0302}$\to$\texttt{0202}.\\[3pt]
        LeetCode's explanation leaves the slice instead (\texttt{1000}, \texttt{1100}, \ldots, \texttt{1202}): also 6 moves.};
\end{tikzpicture}
\end{tracebox}

\begin{tracebox}[{Sample 1: the BFS layers}]
A full BFS from \texttt{0000} (without stopping at the target) gives these layer sizes; 5 codes (the dead ends) are never reached:
\begin{center}
\begin{tikzpicture}[x=0.6cm, y=0.0028cm]
    \foreach \d/\c in {0/1, 1/8, 2/31, 3/86, 4/190, 5/356, 6/576, 7/824, 8/1056, 9/1224, 10/1286,
                       11/1224, 12/1056, 13/824, 14/576, 15/356, 16/192, 17/88, 18/32, 19/8, 20/1} {
        \ifnum\d=6
            \fill[accentorange] (\d-0.35,0) rectangle (\d+0.35,\c);
        \else
            \fill[primaryblue!55] (\d-0.35,0) rectangle (\d+0.35,\c);
        \fi
        \node[font=\tiny, text=darkgray] at (\d,-45) {\d};
    }
    \draw[darkgray] (-0.6,0) -- (20.6,0);
    \foreach \d/\c in {2/31, 6/576, 10/1286, 18/32} \node[font=\tiny, above] at (\d,\c) {\c};
    \node[font=\scriptsize, text=darkgray] at (10,-140) {distance from \texttt{0000}};
\end{tikzpicture}
\end{center}
Layer 2 has 31 codes instead of 32 because \texttt{0101} is dead. The target sits in layer 6 (orange): the early-stopping BFS pops 865 codes before reaching it. The deepest code, \texttt{5555}, is at distance $5\cdot4 = 20$ --- no answer can exceed 20.
\end{tracebox}

\begin{tracebox}[{Sample 3: target walled in}]
\begin{lstlisting}[style=trace]
neighbours of 8888: 9888 7888 8988 8788 8898 8878 8889 8887  -> all dead
one-sided BFS : explores all 9991 live codes reachable from 0000, never sees 8888 -> -1
bidirectional : the target side's first layer is empty -> -1 immediately
\end{lstlisting}
\end{tracebox}

\subsubsection*{6. Test suite}

\begin{testbox}
\centering\small
\begin{tabular}{@{}llcl@{}}
\toprule
\textbf{deadends} & \textbf{target} & \textbf{Expected} & \textbf{Why} \\
\midrule
sample 1 & \texttt{0202} & 6 & detour around dead ends \\
\texttt{8888} & \texttt{0009} & 1 & wrap-around $0 \to 9$ \\
8 neighbours of \texttt{8888} & \texttt{8888} & $-1$ & target isolated \\
\texttt{0000} & \texttt{8888} & $-1$ & start is a dead end \\
\texttt{1234} & \texttt{0000} & 0 & target $=$ start \\
\texttt{9999} & \texttt{5555} & 20 & farthest code \\
8 neighbours of \texttt{0202} & \texttt{0202} & $-1$ & walled target near the start \\
500 random codes & random & --- & stress: BFS $=$ bidirectional \\
\bottomrule
\end{tabular}
\end{testbox}

BFS and bidirectional BFS agree on 4000 random instances (random dead-end sets of up to 500 codes, and dense sets concentrated around the target), and iterative deepening confirms every answer up to 5.

\subsubsection*{7. Common mistakes}

\begin{warningbox}[Mistake 1 --- negative digit on ``down'']
\begin{lstlisting}[style=snippet]
int nd = (g - 1) % 10;             // WRONG: g = 0 gives -1, not 9
int nd = (g + 9) % 10;             // RIGHT
\end{lstlisting}
\end{warningbox}

\begin{warningbox}[Mistake 2 --- forgetting that the start can be dead]
If \texttt{"0000"} is a dead end and the code pushes it anyway, BFS happily walks away from a jammed lock and returns a positive number. Test \texttt{dead[0]} before the loop.
\end{warningbox}

\begin{warningbox}[Mistake 3 --- checking dead ends only for the answer]
Testing ``is the \emph{target} dead?'' is not enough: a dead code in the \emph{middle} of a route jams the lock. Dead codes must never enter the queue.
\end{warningbox}

\begin{warningbox}[Mistake 4 --- bidirectional BFS that returns too early]
Returning at the very first touching edge \emph{in the middle of a layer} can give a non-shortest answer in general graphs. Finish the layer and take the minimum (as above), or test the meeting condition at whole-layer granularity.
\end{warningbox}

\subsubsection*{8. Takeaways}

\begin{itemize}[leftmargin=*]
    \item A combination lock is a 4-dimensional wrap-around grid: $10^4$ vertices, 8 edges each; forbidden codes are deleted vertices.
    \item Encode tuples of digits as integers; the move is a place-value addition.
    \item Without obstacles the distance has a closed formula; obstacles are what make a search necessary.
    \item Bidirectional BFS explores two balls of radius $d/2$ instead of one of radius $d$ --- a big win in fast-growing state spaces.
\end{itemize}

\begin{variantsbox}
\variant{Water jugs}{\varin{https://leetcode.com/problems/water-and-jug-problem/}}
  {instead of 4 wheels, two jugs; moves are fill, empty and pour.}
  {state $(a, b)$ = water in each jug, $(x{+}1)(y{+}1)$ states, encoded as $a(y{+}1)+b$; 6 moves per state.}
  {BFS (or DFS) for reachability of a state with $a + b = T$; a gcd formula also exists (Exercise~E4.3).}
\variant{Operations on one number with a bounded range}{\varin{https://leetcode.com/problems/minimum-operations-to-convert-number/}}
  {one number $x$; moves $x \pm a_i$, $x \oplus a_i$; only values in $[0, 1000]$ can be continued.}
  {vertices $0..1000$; values outside are leaves that only count if they equal the goal.}
  {BFS; test the goal on every generated value before the range check (Exercise~E4.5).}
\variant{More wheels, more digits}{\varnone}
  {$k$ wheels with $b$ digits each.}
  {$b^k$ codes as mixed-radix integers; $2k$ neighbours each.}
  {the same BFS; it is feasible while $b^k \cdot 2k$ fits (e.g.\ $k = 6$, $b = 10$: $1.2\cdot10^7$).}
\variant{Turning two adjacent wheels together}{\varnone}
  {a move may turn wheels $i$ and $i+1$ simultaneously by the same $\pm1$.}
  {same vertices; 6 more edges per code (3 adjacent pairs, 2 directions).}
  {unchanged BFS --- adding a kind of move only adds lines to the neighbour generator.}
\variant{Turns with different costs}{\varnone}
  {turning wheel $i$ costs $c_i$.}
  {same graph, edge weights $c_i$.}
  {BFS no longer works; Dijkstra's algorithm (Chapter 7).}
\end{variantsbox}
